\documentclass[12pt]{article}
\usepackage{amsfonts,amsthm,amsmath,amssymb,mathtools}
\usepackage{mathrsfs}
\usepackage{enumitem}
\usepackage{tikz-cd}
\usepackage{geometry}
\usepackage{mlmodern}
\usepackage{xcolor}

\geometry{letterpaper, portrait, margin=1in}

\setcounter{section}{0}

\renewcommand{\thesection}{\S\arabic{section}}
\renewcommand{\thesubsection}{\arabic{section}}

\DeclareMathOperator{\lcm}{lcm}
\DeclareMathOperator{\im}{im}

\newtheorem{thm}{Theorem}
\newenvironment{theorem}[1]{
	\renewcommand{\thethm}{#1}
	\begin{thm}
		}{\end{thm}}

\newtheorem{lma}{Lemma}
\newenvironment{lemma}[1]{
	\renewcommand{\thelma}{#1}
	\begin{lma}
		}{\end{lma}}

\theoremstyle{definition}
\newtheorem{ex}{}
\newenvironment{exercise}[1]{
	\renewcommand{\theex}{#1}
	\begin{ex}
		}{\end{ex}}

\title{Homework 5}
\author{Connor Mooneyhan}
\date{October 15, 2025}

\begin{document}

\maketitle

Disclaimer: Generative AI was used on problem 3 for understanding the core concepts involved and on problem 6 for help computing the conjugacy classes of $D_8$.

\begin{ex}
	Let $Q_8$ be the Quaternion group.
	Recall that $Q_8 = \left\lbrace  \pm 1, \pm i, \pm j, \pm k \right\rbrace$ with $i^2 = j^2 = k^2 = -1$, $ij = k$, $jk = i$, $ki = j$, $ji = -k$, $kj = -i$, and $ik = -j$.
	\begin{enumerate}[label=(\alph*)]
		\item Find all composition series of $Q_8$ and list the composition factors.
		\item Illustrate Cayley's Theorem applied to $Q_8$ by identifying an isomorphic copy of $Q_8$ as a subgroup of $S_8$.
		\item Prove\footnote{Hint: Prove that if $Q_8$ acts on a set $X$ with $|X| \leq 7$ and $x \in X$, then $\mathrm{Stab}_G(x)$ must contain $\langle -1 \rangle$.} that $Q_8$ is not isomorphic to a subgroup of $S_n$ for any $n < 8$.
	\end{enumerate}
\end{ex}

\begin{ex}
	Compute the orders of the conjugacy classes of $S_6$, and use this to compute the order of the centralizer of each type of permutation in $S_6$.
	Also, use Problem 8 on the previous homework to determine the size of the conjugacy classes of $A_6$.
	Finally, determine the class equation for both $S_6$ and $A_6$.
\end{ex}

\begin{ex}
	Let $G$ be a group of order $3825 = 3^2 \cdot 5^2 \cdot 17$.
	Prove that if $H$ is a normal subgroup of order $17$ in $G$, then $H$ is contained in the center of $G$.
\end{ex}
\begin{proof}
  We let $G$ act on $H$ by conjugation, which induces a homomorphism $\varphi : G \to \mathrm{Aut}(H)$.
  Note that $\mathrm{Aut}(H) = 16$ since $H$ must be cyclic (because 17 is prime).
  The image of this homomorphism must divide both $|G|$ and $|\mathrm{Aut}_H|$, but 3825 and 16 are relatively prime, so $\im \varphi$ must be trivial.
  Thus $H \subset Z(G)$.
\end{proof}

\begin{ex}
	Prove that if $K$ is a characteristic subgroup of $H$ and $H$ is a characteristic subgroup of $G$, then $K$ is a characteristic subgroup of $G$.
	Also show that if $K$ is a characteristic subgroup of $H$ and $H$ is normal in $G$, then $K$ is normal in $G$.
\end{ex}
\begin{proof}
  Given an automorphism $\varphi$ of $G$, $\varphi|H$ is also an automorphism, and so $(\varphi|H)|K = \varphi|K$ is an automorphism.
  Thus $K$ is characteristic in $G$.

  Now consider conjugation by some element $g \in G$.
  This induces an automorphism on $G$, which restricts to an automorphism on $H$ because $H$ is normal.
  Since $K$ is characteristic in $H$, the restriction to $K$ is again an automorphism, and thus $K$ is normal in $G$.
\end{proof}

\begin{ex}
	Let $G$ be a group, and $H$ a proper subgroup of $G$ such that $|G| > |G : H|!$.
	Prove that there must exist a nontrivial normal subgroup of $G$ contained in $H$, so that $G$ is not simple.
\end{ex}

\begin{ex}
	Determine whether or not each of the following class equations are possible for a group of order 10.
	If it is possible, provide a group that realizes it.
	(Here, I am writing the elements of the center as singleton conjugacy classes).
	\begin{enumerate}[label=(\alph*)]
		\item 1 + 1 + 1 + 1 + 1 + 5
		\item 1 + 1 + 1 + 2 + 5
		\item 1 + 2 + 2 + 5
		\item 1 + 2 + 3 + 4
		\item 1 + 1 + 2 + 2 + 2 + 2
	\end{enumerate}
\end{ex}
\begin{proof}
	Let $G$ be the general notation for a candidate group of order 10.
	\begin{enumerate}[label=(\alph*)]
		\item We can see that $|Z(G)| = 5$, and since the only other term is 5, there is one conjugacy class with size 5.
		      Then $|C_G(x)| = 10/5 = 2$, but this is impossible since $Z(G)$ is a subgroup of $C_G(x)$.
		      Thus, this is not the class equation of any group.
		      This reasoning will be abbreviated for future items.
		\item Since $|Z(G)| = 3$ and $|C_G(x)| = 5$ for some $x$, this is not the class equation of a group.
		\item $D_{10}$ satisfies this class equation since the conjugacy classes are $\left\lbrace 1 \right\rbrace$, $\left\lbrace r, r^4 \right\rbrace$, $\left\lbrace r^2, r^3 \right\rbrace$, $\left\lbrace s, rs, r^2s, r^3s, r^4s \right\rbrace$.
		\item Since 3 does not divide 10, this is not the class equation of a group.
		\item Since $|Z(G)| = 2$ and $|C_G(x)| = 5$ for some $x$, this is not the class equation of a group.
	\end{enumerate}
\end{proof}

\begin{ex}
	Suppose that $G$ is a group with class equation $1 + 4 + 5 + 5 + 5$.
	\begin{enumerate}[label=(\alph*)]
		\item Does $G$ have a subgroup of order 5? If so, is it a normal subgroup?
		\item Does $G$ have a subgroup of order 4? If so, is it a normal subgroup?
	\end{enumerate}
\end{ex}
\begin{proof}
  Let $x$ be an element of the conjugacy class of order 4, and let $y$ be an element of one of the conjugacy classes of order 5.
  \begin{enumerate}[label=(\alph*)]
    \item Yes, $C_G(x)$ is such a subgroup, since $|C_G(x)| = |G|/4 = 20/4 = 5$.
      The union of conjugacy classes is normal if it is a subgroup, and indeed we have a subgroup in the union of the trivial conjugacy class and the order-4 conjugacy class, so we have a normal subgroup of order 5.
    \item Yes, $C_G(y)$ is such a subgroup, since $|C_G(y)| = |G|/5 = 20/5 = 4$.
      Since 4 cannot be written as a combination of 1 with any other of the orders of conjugacy classes, it cannot be a normal subgroup.
  \end{enumerate}
\end{proof}

\begin{ex}
	Let $G$ be a nonabelian group of order $p^3$ with $p$ a prime integer.
	Determine $|Z(G)|$.
	Next, determine $|C_G(x)|$ for $x \notin Z(G)$.
	Finally, determine the class equation of $G$ using this information.
\end{ex}

\end{document}
